% Condensed Introduction. Verbatim pulls from sources/introduction.tex and
% sources/contributions.tex, cut to: thesis statement -> historical pivot
% (knowledge came from the system's own records; results did not transfer) ->
% the question and its three sub-questions -> the seven contributions -> roadmap.
% No figures (budget: 0 for the introduction).

\partepigraph{``The biggest lesson that can be read from 70 years of AI research is
that general methods that leverage computation are ultimately the most effective, and
by a large margin.''}{Richard S. Sutton, \textit{The Bitter Lesson}}

\noindent How can we teach an artificial intelligence to read hospital notes when hospital
notes cannot be used to teach it? The notes themselves would be
the best examples because they contain the words, shortcuts, and clinical habits
the system needs to learn. But they describe real patients, so they stay inside
the hospital, and a system trained on them may have to stay there too. This
thesis asks whether public biomedical text can be used instead.

The question is old. In 1959, Robert Ledley and Lee Lusted set a medical goal for
the new field of artificial intelligence: a computer that could weigh a patient's
signs toward a likely disease~\citep{ledley_lusted_1959}.
Their question is also
the question of this thesis: where does the knowledge that drives such a system
come from, and can a machine be given enough of it to be useful in care? The early
medical systems answered from their own records. De~Dombal's acute-abdomen system
reached 91.8\% diagnostic accuracy on 304 patients at a single centre in
1972~\citep{dedombal_1972}, but when it was taken to a multicentre trial across
eight centres and 16{,}737 patients, accuracy with the aid in place rose only to
65.3\%~\citep{adams_1986}. A system tuned on Leeds patients did not carry its
91.8\% to other hospitals, because the probabilities it relied on came from one
place and the new patients came from many. A learning system depends on the data
its knowledge was built from, so a result measured in one setting may not hold in
another.

In the 1970s, medical AI shifted toward \emph{expert systems}, which stored
specialist knowledge as explicit rules. MYCIN, at Stanford, advised on antibiotics
through IF--THEN rules~\citep{mycin_1975}; in a blind test, eight infectious-disease
specialists judged it acceptable in 65\% of cases, against a faculty mean of
55.5\%~\citep{mycin_eval_1979}. It passed the benchmark and was never used to treat a
patient~\citep{internist1_1982}. A program could win a blind benchmark and still be
unusable, because winning the benchmark was not what physicians were asking for.

Machine learning later changed how systems acquired knowledge. Instead of being
told the rules, systems learned what to do from large sets of examples, and the
gains were large: by 2016 a network trained on 128{,}175 retinal photographs read
them for diabetic eye disease about as accurately as
ophthalmologists~\citep{gulshan_2016}. In the biomedical literature, expert-system
papers peak around 1991, then a deep-learning line near zero before 2015 climbs past
twenty thousand papers a year. These results raised a new requirement, the one that
organises this thesis: a system that learns from examples needs many well-matched
examples, and clinical text does not circulate freely.

Most modern natural language processing starts by pretraining a language model on
large collections of raw text, before adapting it to a task. In the clinical
setting this path is hard to follow. Real clinical notes are personal data and
cannot be shared, so France has no public corpus of them to pretrain on. A model
trained inside one hospital on its own records carries the same constraint and
cannot be released to others. A shareable clinical model has to be built from
public text.

This thesis asks one question: can public data support language models for French
clinical information extraction, without training on sensitive clinical records?
The manuscript answers it through three questions.

\begin{itemize}
  \item Can we assemble a useful public French biomedical corpus, and does its content and quality change what the model learns?
  \item Given such a corpus, how should we pretrain a language model on it: by adapting an existing model, choosing the training objective, or adding knowledge from outside the text?
  \item How can we adapt these models to clinical extraction tasks when labelled clinical data is scarce, through model architecture and synthetic data?
\end{itemize}

\noindent This thesis makes seven contributions.

\begin{itemize}
  \item \emph{biomed-fr}, a public French biomedical corpus of 413 million words, used to train CamemBERT-bio. The adapted model finds biomedical terms in French text, such as drugs and diseases, more accurately than the general model it started from, with a 2.54-point average gain and at a fraction of the cost of training from scratch~\citep{touchent2023camembertbio}.
  \item \emph{Biomed-Enriched}, a labelling of the open biomedical literature (PubMed Central) by content type and quality. It identifies about two million paragraphs of clinical cases and improves a model trained further on them~\citep{touchent-etal-2026-biomed}.
  \item \emph{MC-Bio}, a ten-billion-token French biomedical corpus, and a study of which content types and quality signals help. The experiments show that some common filters degrade performance.
  \item A pretraining method that trains the model on a different task for a while and then switches back to the main objective, improving the model at no extra cost.
  \item \emph{OntoBook}, a way to add structured medical knowledge to pretraining by turning a medical knowledge base into training text.
  \item Model designs, MC-bio-gliner and MC-bio-embed, that find clinical terms from only a handful of labelled examples.
  \item A method for generating synthetic clinical reports by rewriting public biomedical text into hospital style, so that extractors can be trained without access to real hospital data.
\end{itemize}

\noindent The parts follow these questions. The first asks what biomedical text is and how
far its public forms are from the clinical notes we care about. The second builds
and studies a public French biomedical corpus. The third pretrains language models
on it, through model adaptation, training objectives, and knowledge enrichment. The
final part adapts these models to clinical extraction tasks under scarce labelled
data. The conclusion then opens onto a question this work leaves behind: once such
models are built, how they behave when someone actually uses them becomes its own
problem.
